\section{Introduction}
\label{sec:intro}

Modern enterprise campuses and data-center fabrics are operated under continuous change: VLAN and ACL updates, EVPN/VXLAN overlays, QoS policy edits, and firmware-driven control-plane behavior. Cognitive NetOps and AIOps promise earlier anomaly detection, failure-horizon prediction, explainable RCA, service-impact estimation, and remediation recommendations~\cite{dang2019aiops,notaro2021survey,almasan2022ndt}. Beyond purely technical detectors, enterprise modernization also depends on governance maturity, workforce adaptation, and cybersecurity readiness when operators adopt digital twins and automation~\cite{rah2026dtbeyond,rah2026dtfails,rah2026govresistance,rah2026govdtcyber}. In practice, progress is slowed by fragmented public resources: topology archives without multi-modal telemetry~\cite{knight2011topologyzoo}, security flow corpora that omit enterprise L2/L3 configuration semantics~\cite{moustafa2015unsw}, and telemetry/RCA artifacts that do not jointly expose recovery and impact labels under a temporal leakage protocol.

This paper addresses that gap with two coupled contributions. First, \textbf{ECNetBench} is a synthetic enterprise cognitive-networking benchmark whose schema, failure taxonomy, and label definitions are designed for leakage-safe supervised learning under AOS-CX-style operational assumptions (inventory, VLAN/ACL/QoS, routing, counters, syslog/alerts, NAE-like analytics, and recovery/impact fields). Second, we implement and evaluate an \textbf{Enterprise Cognitive Network (ECN-v3)}: a Digital Twin graph plus a multi-agent stack (Perception, Anomaly, Prediction, RCA, Impact, Healing) with leakage-safe feature enrichment, \emph{anchored} telemetry-first fusion, and TreeSHAP-explained RCA. The novelty is the coupled benchmark, explicit temporal leakage protocol, and evidence-selected fusion evaluation on six frozen instances---not a claim of absolute historical priority over all related AIOps or digital-twin systems.

Our contributions are: (i)~ECNetBench design (relational/graph schema, causal incident modeling, leakage-safe labels, and six frozen instances); (ii)~the ECN-v3 framework with architecture isolation of anchored fusion versus stacking; (iii)~multi-seed evaluation with baselines, ablations, calibration, confidence intervals, and paired significance tests under $n{=}6$; and (iv)~a reproducible package with checksums and relative-path evaluation. Empirically, the final anchored T1 head improves mean AUPRC relative to the prior v2 configuration and strong tabular baselines, while paired tests versus random forest remain non-significant and stacking is a negative ablation; detailed metrics appear in Section~\ref{sec:results}.

The overall architecture is introduced in Fig.~\ref{fig:architecture} (Section~\ref{sec:framework}). Section~\ref{sec:related} situates related work; Sections~\ref{sec:bench}--\ref{sec:framework} detail the benchmark and framework; Sections~\ref{sec:eval}--\ref{sec:results} present protocol and results; Sections~\ref{sec:disc}--\ref{sec:conc} discuss implications, limitations, and conclusions.
